\begin{abstract}
As decentralized finance protocols have grown to secure hundreds of billions of dollars in on-chain value, so too has the cost of failure: despite tens of millions spent annually on third-party audits, formal verification, and bug bounty programs, smart contract vulnerabilities continue to cause billions of dollars in losses. Among these, reentrancy -- where an external call re-enters the calling contract before its state updates complete -- remains one of the most common and damaging vulnerability classes. Because it is a purely structural property of the call graph rather than a semantic bug, it is uniquely suited to prevention through language design. We present a Rust-inspired smart contract language targeting the EVM that introduces two key safety mechanisms:
1. A compile-time effect system requiring developers to explicitly annotate state modifications.
2. Automatically generated reentrancy guards, with declarative annotations that allow developers to selectively re-enable specific reentrant call paths when required.
Notably, the guarantees provided by the reentrancy guards allow the compiler to propagate side effect constraints even across reentrant call paths, which is not possible in any other language without sophisticated static analysis tools. The explicit effect annotations make storage mutation patterns machine-checked and visible at function boundaries, making smart contract code simpler to secure for developers, auditors, and bug bounty hunters. Together, these mechanisms render reentrancy vulnerabilities virtually impossible. We evaluate the practical impact of these features against real-world DeFi exploits (2021-2024) and find that 53/54 reentrancy incidents would have been prevented, the sole outlier being a compiler bug rather than a developer oversight.
\end{abstract}